%----------------------------------------------------------------------------------------
% GONIOMETRISCHE SUBSTITUTIES. Alle opgaven zijn eigen werk.
%
% Opbouw: Probleem, Start, leerstof, en op een nieuwe pagina de Opdrachten, Koppel en Test jezelf.
% Labels van dit hoofdstuk beginnen met h2-.
%----------------------------------------------------------------------------------------

\mychapter{2}{Goniometrische substituties}

\begin{probleem}
\naast{In de kelder ligt een cilindervormige stookolietank: 3~m lang, met een diameter van 2~m. Om te weten hoeveel olie er nog in zit, steek je een peilstok verticaal tot op de bodem. De olie komt tot 50~cm hoog: een kwart van de hoogte. Zit de tank dan voor een kwart vol? Hoeveel liter is het? Het volume is de lengte maal de oppervlakte van het blauwe stuk van de cirkel. Die oppervlakte is een integraal van $2\sqrt{1 - y^2}$, en voor zo'n wortel kende je tot nu toe geen primitieve.}{\begin{tikzpicture}[scale=1.05,font=\scriptsize]
  \fill[accent!35] (-.866,-.5) arc[start angle=210,end angle=330,radius=1] -- cycle;
  \draw[accent,thick] (0,0) circle (1);
  \draw[accent] (-.866,-.5) -- (.866,-.5);
  \draw[->,grijs] (-1.3,0) -- (1.4,0) node[right,text=black] {$x$};
  \draw[->,grijs] (0,-1.25) -- (0,1.35) node[above,text=black] {$y$};
  \draw[fel!80!black,thick] (.35,-1) -- (.35,1.15);
  \draw[lijngrijs,|-|] (1.25,-1) -- (1.25,-.5) node[midway,right] {0,5 m};
  \node[below left] at (0,-1) {$-1$};
  \node[above left] at (0,1) {1};
\end{tikzpicture}}
\end{probleem}

\lkr{Leerplan Complementair (III-WiCo-d): LPD 18. Voorkennis uit pakket 4: de substitutieregel, partiële integratie en de standaardvormen van Integralen (hoofdstuk 3), de formules van Carnot en $1 + \tan^2 x = \tfrac{1}{\cos^2 x}$ (Goniometrische functies), en de afgeleiden van $\tan$ en $\bgsin$. Uit hoofdstuk 1: partieelbreuken, voor $\int \tfrac{1}{\cos x}\dx$ en na de t-formules. De t-formules kwamen in Integralen al voor in een opdracht van niveau 3 (de Weierstrasssubstitutie): hier worden ze leerstof, met een meetkundige betekenis. Hyperbolische functies zijn niet nodig; ze staan alleen in een opdracht van Integralen.}

In de cursus Integralen stelde je altijd $u = g(x)$: je gaf een stuk van de integrand een nieuwe naam. Bij een wortel zoals $\sqrt{1 - y^2}$ uit het probleem lukt dat niet, want de afgeleide van $1 - y^2$ staat er niet naast. In dit hoofdstuk draaien we de substitutie om: we vervangen de veranderlijke zelf door een goniometrische functie. Waarom zou dat helpen? Omdat $1 - \sin^2 u = \cos^2 u$, en uit een kwadraat kun je de wortel trekken. Zo verdwijnen wortels van $a^2 - x^2$, $a^2 + x^2$ en $x^2 - a^2$. Op het einde zie je nog een substitutie die de goniometrie zelf laat verdwijnen: de t-formules.

\begin{oefeningen}[Start]
\oef{h2-start-1}

\oef{h2-verken}
\end{oefeningen}

\lkr{Laat oefening \ref{oef:h2-verken} maken vóór de uitleg. Leerlingen kennen het antwoord al uit de meetkunde, zien dat de gewone substitutie faalt, en rekenen dan zelf de nieuwe substitutie uit. Vraag c is het hele idee van het hoofdstuk.}

%========================================================================================
\section{Substitutie in de andere richting}

In de Verken-opgave vervang je $x$ door $\sin u$. Dat is een substitutie in de omgekeerde richting: niet $u = g(x)$, maar $x = g(u)$. Mag dat zomaar? Ja: het is dezelfde substitutieregel, van rechts naar links gelezen. Let wel op de grenzen: bij een bepaalde integraal moet je voor elke grens van $x$ een waarde van $u$ vinden.

\begin{theorie}
\begin{stelling}{Substitutie $x = g(u)$}
Is $f$ continu, en heeft $g$ een continue afgeleide op $[\alpha, \beta]$ met $g(\alpha) = a$ en $g(\beta) = b$, dan is

\vspace{1mm}
\qquad $\formule{\int_a^b f(x)\dx = \int_\alpha^\beta f(g(u))\,g'(u)\du}$
\end{stelling}

\begin{bewijs}
Neem een primitieve $F$ van $f$. Volgens de kettingregel is
\bewijsregel{\bigl(F(g(u))\bigr)' = f(g(u))\,g'(u) ,}
dus $F \circ g$ is een primitieve van de integrand rechts. De hoofdstelling geeft
\bewijsregel{\ig_\alpha^\beta f(g(u))\,g'(u)\du = F(g(\beta)) - F(g(\alpha)) = F(b) - F(a) = \ig_a^b f(x)\dx .}
\end{bewijs}
\end{theorie}

In de notatie van Leibniz schrijf je $x = g(u)$ en $\mathrm{d}x = g'(u)\du$, net zoals in Integralen. Bij een onbepaalde integraal moet je op het einde terug naar $x$. Dan moet je $u$ uit $x = g(u)$ kunnen halen: kies $u$ dus op een interval waar $g$ een inverse heeft. Voor $x = \sin u$ is dat $\bigl[-\tfrac{\pi}{2}, \tfrac{\pi}{2}\bigr]$, met $u = \bgsin x$.

\voorbeeld{}{Bereken $\ig_0^2 \sqrt{4 - x^2}\dx$, en controleer met een figuur.\label{vb:h2-andere}}{%
\uitleg{Stel $x = 2\sin u$: dan wordt $4 - x^2 = 4\cos^2 u$, en verdwijnt de wortel. Voor de grenzen zoeken we de $u$ waarbij $x = 0$ en $x = 2$.}
\naast{\begin{stappenvrij}{6.8cm}{3.4cm}
\text{Stel } x = 2\sin u, \quad \mathrm{d}x = 2\cos u\du & $x = 0$ bij $u = 0$, en $x = 2$ bij $u = \tfrac{\pi}{2}$. \\[1ex]
\sqrt{4 - 4\sin^2 u} = 2\sqrt{\cos^2 u} = 2\cos u & Op $\bigl[0, \tfrac{\pi}{2}\bigr]$ is $\cos u \geqslant 0$. \\[1ex]
4\ig_0^{\pi/2} \cos^2 u\du & Vul alles in: $2\cos u \cdot 2\cos u\du$. \\[2.5ex]
= 4 \grens{\dfrac u2 + \dfrac{\sin 2u}{4}}{0}{\pi/2} = \pi & Met de formule van Carnot.
\end{stappenvrij}}{\begin{tikzpicture}[scale=.75,font=\scriptsize]
  \fill[accent!30] (0,0) -- (2,0) arc[start angle=0,end angle=90,radius=2] -- cycle;
  \rooster{-2}{1}{2}{0}{1}{2}
  \draw[accent,thick] (2,0) arc[start angle=0,end angle=180,radius=2];
\end{tikzpicture}}
De grafiek van $\sqrt{4 - x^2}$ is een halve cirkel met straal 2, dus de integraal is een kwart van de schijf: $\tfrac14 \cdot \pi \cdot 2^2 = \pi$.}

\begin{werkwijze}
Stel $x = g(u)$ en schrijf $\mathrm{d}x = g'(u)\du$. Zoek bij elke grens van $x$ een grens van $u$, en schrijf de hele integrand in $u$. Bij een onbepaalde integraal ga je op het einde terug naar $x$.
\end{werkwijze}

\nujij{Bereken $\int_0^1 \dfrac{x^2}{\sqrt{1 - x^2}}\dx$ met $x = \sin u$.}{Grenzen $0$ en $\tfrac{\pi}{2}$: $\int_0^{\pi/2} \sin^2 u\du = \tfrac{\pi}{4}$.}

\begin{oefeningen}
\oef{h2-andere}
\end{oefeningen}

%========================================================================================
\section{Een wortel van \texorpdfstring{$a^2 - x^2$}{a²-x²}}

\subsection{Met de sinus}

Hoe vind je een primitieve van $\sqrt{9 - x^2}$, zonder grenzen? De substitutie $x = 3\sin u$ maakt van de wortel $3\cos u$. Nieuw is de weg terug: op het einde staan er $\sin u$, $\cos u$ en $u$, en die moet je in $x$ uitdrukken. Een rechthoekige driehoek helpt. Uit $\sin u = \tfrac x3$ volgt een driehoek met schuine zijde 3 en overstaande zijde $x$, en Pythagoras geeft de aanliggende zijde.

\voorbeeld{}{Bereken $\ig \sqrt{9 - x^2}\dx$.\label{vb:h2-sin}}{%
\uitleg{Stel $x = 3\sin u$ met $u \in \bigl[-\tfrac{\pi}{2}, \tfrac{\pi}{2}\bigr]$. Dan is $\cos u \geqslant 0$, en geeft de wortel $3\cos u$.}
\naast{\begin{stappenvrij}{7.6cm}{3.6cm}
\text{Stel } x = 3\sin u, \quad \mathrm{d}x = 3\cos u\du & En $\sqrt{9 - 9\sin^2 u} = 3\cos u$. \\[1ex]
\ig 3\cos u \cdot 3\cos u\du = 9\ig \cos^2 u\du & Vul alles in. \\[2.5ex]
= \dfrac92\bigl(u + \sin u\cos u\bigr) + C & Carnot, en $\tfrac{\sin 2u}{2} = \sin u\cos u$. \\[1ex]
= \dfrac92 \bgsin\dfrac x3 + \dfrac x2\sqrt{9 - x^2} + C & Lees $\sin u = \tfrac x3$ en $\cos u = \tfrac{\sqrt{9 - x^2}}{3}$ af in de driehoek.
\end{stappenvrij}}{\begin{tikzpicture}[scale=.75,font=\scriptsize]
  \draw[accent,thick] (0,0) -- (2.6,0) -- (2.6,1.5) -- cycle;
  \draw[lijngrijs] (2.4,0) -- (2.4,.2) -- (2.6,.2);
  \draw (.65,0) arc[start angle=0,end angle=30,radius=.65];
  \node at (.95,.22) {$u$};
  \node[above left] at (1.3,.75) {$3$};
  \node[right] at (2.6,.75) {$x$};
  \node[below] at (1.3,0) {$\sqrt{9 - x^2}$};
\end{tikzpicture}}}

Met $3$ vervangen door een willekeurige $a > 0$ vind je net zo een formule. Die controleren we met de afgeleide, en daarna zien we wat ze betekent.

\begin{theorie}
\begin{stelling}{Primitieve van $\sqrt{a^2 - x^2}$}
Voor $a > 0$ en $x \in [-a, a]$ is \qquad $\formule{\int \sqrt{a^2 - x^2}\dx = \dfrac{a^2}{2}\bgsin\dfrac xa + \dfrac x2\sqrt{a^2 - x^2} + C}$
\end{stelling}

\begin{bewijs}
Leid het rechterlid af op $]-a, a[$, met $(\bgsin\tfrac xa)' = \tfrac{1}{\sqrt{a^2 - x^2}}$:
\bewijsregel{\dfrac{a^2}{2\sqrt{a^2 - x^2}} + \dfrac12\sqrt{a^2 - x^2} - \dfrac{x^2}{2\sqrt{a^2 - x^2}} = \dfrac{a^2 + (a^2 - x^2) - x^2}{2\sqrt{a^2 - x^2}} = \sqrt{a^2 - x^2} .}
\end{bewijs}
\end{theorie}

\naast{De twee termen hebben een meetkundige betekenis. De oppervlakte onder de cirkel $y = \sqrt{a^2 - x^2}$ van 0 tot $x$ bestaat uit een driehoek en een cirkelsector. De driehoek heeft basis $x$ en hoogte $\sqrt{a^2 - x^2}$: dat is de term $\tfrac x2\sqrt{a^2 - x^2}$. De sector heeft straal $a$ en middelpuntshoek $u = \bgsin\tfrac xa$, en dus oppervlakte $\tfrac{a^2}{2}\,u$. De substitutie $x = a\sin u$ rekent dus in feite met die hoek.}{\begin{tikzpicture}[scale=1.75,font=\scriptsize]
  \fill[accent!35] (0,0) -- (0,1.4) arc[start angle=90,end angle=55,radius=1.4] -- cycle;
  \fill[accent!15] (0,0) -- (.803,1.147) -- (.803,0) -- cycle;
  \draw[->,grijs] (-.2,0) -- (1.7,0) node[right,text=black] {$x$};
  \draw[->,grijs] (0,-.1) -- (0,1.65) node[above,text=black] {$y$};
  \draw[accent,thick] (1.4,0) arc[start angle=0,end angle=105,radius=1.4];
  \draw[accent] (0,0) -- (.803,1.147) -- (.803,0);
  \draw (0,.35) arc[start angle=90,end angle=55,radius=.35];
  \node at (.1,.5) {$u$};
  \node[below] at (.803,0) {$x$};
  \node[below] at (1.4,0) {$a$};
  \node at (.58,.25) {driehoek};
  \node at (.3,1.05) {sector};
\end{tikzpicture}}

\needspace{10\baselineskip}
Ook als de wortel in de noemer staat, of met een macht, werkt dezelfde substitutie.

\voorbeeld{}{Bereken $\ig \dfrac{1}{(4 - x^2)^{3/2}}\dx$.\label{vb:h2-sin2}}{%
\uitleg{De noemer is $\bigl(\sqrt{4 - x^2}\bigr)^3$: stel $x = 2\sin u$, dan wordt dat $8\cos^3 u$.}
\begin{stappenbreed}
\text{Stel } x = 2\sin u, \quad \mathrm{d}x = 2\cos u\du & Met $u \in \bigl]-\tfrac{\pi}{2}, \tfrac{\pi}{2}\bigr[$. \\[1ex]
\ig \dfrac{2\cos u}{8\cos^3 u}\du = \dfrac14\ig \dfrac{1}{\cos^2 u}\du = \dfrac14\tan u + C & Een basisintegraal. \\[2.5ex]
= \dfrac{x}{4\sqrt{4 - x^2}} + C & In de driehoek: overstaand $x$, aanliggend $\sqrt{4 - x^2}$.
\end{stappenbreed}}

\begin{werkwijze}
Bij een wortel van $a^2 - x^2$: stel $x = a\sin u$ met $u \in \bigl[-\tfrac{\pi}{2}, \tfrac{\pi}{2}\bigr]$, dan is $\sqrt{a^2 - x^2} = a\cos u$ en $\mathrm{d}x = a\cos u\du$. Om terug te keren naar $x$ teken je een rechthoekige driehoek met schuine zijde $a$, overstaande zijde $x$ en aanliggende zijde $\sqrt{a^2 - x^2}$, en $u = \bgsin\tfrac xa$.
\end{werkwijze}

\nujij{Bereken $\int \dfrac{1}{(1 - x^2)^{3/2}}\dx$.}{$x = \sin u$: $\int \tfrac{1}{\cos^2 u}\du = \tan u + C = \dfrac{x}{\sqrt{1 - x^2}} + C$.}

\begin{oefeningen}
\oef{h2-sin}
\end{oefeningen}

\subsection{Terug naar het probleem}

\voorbeeld{}{Terug naar het probleem. In de kelder ligt een cilindervormige stookolietank: 3~m lang, met een diameter van 2~m. Om te weten hoeveel olie er nog in zit, steek je een peilstok verticaal tot op de bodem. De olie komt tot 50~cm hoog: een kwart van de hoogte. Zit de tank dan voor een kwart vol? Hoeveel liter is het?\label{vb:h2-terug}}{%
\naast{\uitleg{Het volume is de lengte maal de oppervlakte van de doorsnede die met olie gevuld is: het blauwe stuk van de cirkel in de figuur. Leg de oorsprong in het midden van de cirkel, met eenheid 1~m. De olie vult de cirkel van $y = -1$ tot $y = -\tfrac12$. Snij dat stuk in dunne horizontale stroken: op hoogte $y$ is een strook $2\sqrt{1 - y^2}$ breed. Dat is de lengte van de koorde, want $x^2 + y^2 = 1$.}}{\begin{tikzpicture}[scale=1.05,font=\scriptsize]
  \fill[accent!35] (-.866,-.5) arc[start angle=210,end angle=330,radius=1] -- cycle;
  \draw[accent,thick] (0,0) circle (1);
  \draw[accent] (-.866,-.5) -- (.866,-.5);
  \draw[->,grijs] (-1.3,0) -- (1.4,0) node[right,text=black] {$x$};
  \draw[->,grijs] (0,-1.25) -- (0,1.35) node[above,text=black] {$y$};
  \draw[fel!80!black,thick] (.35,-1) -- (.35,1.15);
  \draw[lijngrijs,|-|] (1.25,-1) -- (1.25,-.5) node[midway,right] {0,5 m};
  \node[below left] at (0,-1) {$-1$};
  \node[above left] at (0,1) {1};
\end{tikzpicture}}
\begin{stappenbreed}
A = \ig_{-1}^{-1/2} 2\sqrt{1 - y^2}\,\mathrm{d}y & De oppervlakte van het blauwe stuk. \\[3ex]
= \Bigl[\bgsin y + y\sqrt{1 - y^2}\,\Bigr]_{-1}^{-1/2} & De formule met $a = 1$, maal 2. \\[3ex]
= \Bigl(-\dfrac{\pi}{6} - \dfrac12 \cdot \dfrac{\sqrt3}{2}\Bigr) - \Bigl(-\dfrac{\pi}{2} - 0\Bigr) = \dfrac{\pi}{3} - \dfrac{\sqrt3}{4} \approx 0{,}614 & $\bgsin\bigl(-\tfrac12\bigr) = -\tfrac{\pi}{6}$ en $\bgsin(-1) = -\tfrac{\pi}{2}$. \\[3ex]
V = 3A \approx 1{,}843 \text{ m}^3 & Het volume is de lengte maal de doorsnede. \\[1ex]
\dfrac{V}{3\pi} \approx 0{,}196 & De volle tank heeft $\pi \cdot 1^2 \cdot 3 = 3\pi \approx 9{,}42$~m$^3$.
\end{stappenbreed}
Er zit ongeveer 1840 liter in, minder dan 20\,\% van de tank, en niet 25\,\%: onderaan is de tank smal. Dat zie je ook in de formule. De sector heeft hoek $\tfrac{2\pi}{3}$, een derde van de cirkel, maar daar gaat de driehoek met top in het middelpunt nog af.}

%========================================================================================
\section{Een wortel van \texorpdfstring{$a^2 + x^2$}{a²+x²}}

\subsection{Met de tangens}

Bij $a^2 + x^2$ helpt de sinus niet, want $a^2 + a^2\sin^2 u$ is geen kwadraat. Welke goniometrische formule maakt van een som een kwadraat? Die met de tangens: $1 + \tan^2 u = \dfrac{1}{\cos^2 u}$. Stel dus $x = a\tan u$ met $u \in \bigl]-\tfrac{\pi}{2}, \tfrac{\pi}{2}\bigr[$. Dan is $a^2 + x^2 = \dfrac{a^2}{\cos^2 u}$, en omdat $\cos u > 0$ is de wortel $\dfrac{a}{\cos u}$.

\voorbeeld{}{Bereken $\ig \dfrac{1}{(x^2 + 4)^{3/2}}\dx$.\label{vb:h2-tan}}{%
\uitleg{Stel $x = 2\tan u$. De afgeleide van $\tan u$ is $\tfrac{1}{\cos^2 u}$. Voor de terugweg teken je een driehoek met overstaande zijde $x$ en aanliggende zijde 2.}
\naast{\begin{stappenvrij}{7.4cm}{3.6cm}
\text{Stel } x = 2\tan u, \quad \mathrm{d}x = \dfrac{2}{\cos^2 u}\du & En $(x^2 + 4)^{3/2} = \dfrac{8}{\cos^3 u}$. \\[2.5ex]
\ig \dfrac{2}{\cos^2 u} \cdot \dfrac{\cos^3 u}{8}\du = \dfrac14\ig \cos u\du & Delen door een breuk is vermenigvuldigen met het omgekeerde. \\[2.5ex]
= \dfrac14\sin u + C = \dfrac{x}{4\sqrt{x^2 + 4}} + C & In de driehoek is de schuine zijde $\sqrt{x^2 + 4}$.
\end{stappenvrij}}{\begin{tikzpicture}[scale=.75,font=\scriptsize]
  \draw[accent,thick] (0,0) -- (2.2,0) -- (2.2,1.6) -- cycle;
  \draw[lijngrijs] (2,0) -- (2,.2) -- (2.2,.2);
  \draw (.6,0) arc[start angle=0,end angle=36,radius=.6];
  \node at (.9,.22) {$u$};
  \node[above left] at (1.1,.8) {$\sqrt{x^2 + 4}$};
  \node[right] at (2.2,.8) {$x$};
  \node[below] at (1.1,0) {$2$};
\end{tikzpicture}}}

\voorbeeld{}{Bereken $\ig_0^1 \dfrac{1}{(1 + x^2)^2}\dx$.\label{vb:h2-tan2}}{%
\uitleg{In hoofdstuk 1 lukte dit alleen met een omweg via partiële integratie. Met $x = \tan u$ gaat het rechtstreeks.}
\begin{stappen}
\text{Stel } x = \tan u, \quad \mathrm{d}x = \dfrac{1}{\cos^2 u}\du & Grenzen: $u = 0$ en $u = \tfrac{\pi}{4}$. \\[2ex]
\ig_0^{\pi/4} \cos^4 u \cdot \dfrac{1}{\cos^2 u}\du = \ig_0^{\pi/4} \cos^2 u\du & Want $\dfrac{1}{(1 + \tan^2 u)^2} = \cos^4 u$. \\[2.5ex]
= \grens{\dfrac u2 + \dfrac{\sin 2u}{4}}{0}{\pi/4} = \dfrac{\pi}{8} + \dfrac14 & Carnot.
\end{stappen}}

\subsection{De integraal van \texorpdfstring{$\frac{1}{\cos x}$}{1/cos x}}

Wat geeft $x = \tan u$ bij $\int \dfrac{1}{\sqrt{x^2 + 1}}\dx$? De wortel wordt $\dfrac{1}{\cos u}$ en $\mathrm{d}x = \dfrac{\mathrm{d}u}{\cos^2 u}$, dus de integrand wordt $\dfrac{1}{\cos u}$. Die staat niet in de tabel. Hier helpt hoofdstuk 1: vermenigvuldig teller en noemer met $\cos u$, en er staat een rationale functie van $\sin u$.

\begin{theorie}
\begin{stelling}{Primitieve van $\frac{1}{\cos x}$}
Op een interval waar $\cos x \neq 0$ is \qquad $\formule{\int \dfrac{1}{\cos x}\dx = \ln\Bigl|\dfrac{1 + \sin x}{\cos x}\Bigr| + C}$
\end{stelling}

\begin{bewijs}
Schrijf $\dfrac{1}{\cos x} = \dfrac{\cos x}{\cos^2 x} = \dfrac{\cos x}{1 - \sin^2 x}$, en stel $s = \sin x$, $\mathrm{d}s = \cos x\dx$. Met partieelbreuken is
\bewijsregel{\ig \dfrac{\mathrm{d}s}{1 - s^2} = \dfrac12\ig \Bigl(\dfrac{1}{1 + s} + \dfrac{1}{1 - s}\Bigr)\mathrm{d}s = \dfrac12\ln\dfrac{1 + s}{1 - s} + C ,}
want $-1 < s < 1$. Vermenigvuldig teller en noemer van de breuk met $1 + s$:
\bewijsregel{\dfrac12\ln\dfrac{(1 + \sin x)^2}{1 - \sin^2 x} = \dfrac12\ln\dfrac{(1 + \sin x)^2}{\cos^2 x} = \ln\Bigl|\dfrac{1 + \sin x}{\cos x}\Bigr| .}
\end{bewijs}
\end{theorie}

Deze integraal heeft een merkwaardige geschiedenis. Op een zeekaart volgens Mercator (1569) staan de breedtecirkels steeds verder uit elkaar naarmate je naar de pool gaat, zodat een vaste kompaskoers een rechte lijn wordt. De afstand van de evenaar tot breedte $\varphi$ is op die kaart een integraal van $\tfrac{1}{\cos x}$ van 0 tot $\varphi$. Zeelieden rekenden ze eerst met tabellen benaderend uit. Pas in de 17de eeuw zag men dat er een logaritme uitkomt.

Nu kunnen we de tangenssubstitutie afmaken voor de eenvoudigste wortel met een som.

\begin{theorie}
\begin{stelling}{Primitieve van $\frac{1}{\sqrt{x^2 + a^2}}$}
Voor $a > 0$ is \qquad $\formule{\int \dfrac{1}{\sqrt{x^2 + a^2}}\dx = \ln\bigl(x + \sqrt{x^2 + a^2}\,\bigr) + C}$
\end{stelling}

\begin{bewijs}
Stel $x = a\tan u$ met $u \in \bigl]-\tfrac{\pi}{2}, \tfrac{\pi}{2}\bigr[$. Dan is $\sqrt{x^2 + a^2} = \dfrac{a}{\cos u}$ en $\mathrm{d}x = \dfrac{a}{\cos^2 u}\du$, dus
\bewijsregel{\ig \dfrac{\cos u}{a} \cdot \dfrac{a}{\cos^2 u}\du = \ig \dfrac{1}{\cos u}\du = \ln\Bigl(\dfrac{1}{\cos u} + \tan u\Bigr) + C ,}
en dat argument is positief. In $x$ is het $\ln\dfrac{\sqrt{x^2 + a^2} + x}{a} = \ln\bigl(x + \sqrt{x^2 + a^2}\bigr) - \ln a$, en $-\ln a$ zit in de constante.
\end{bewijs}
\end{theorie}

\subsection{De lengte van een parabool}

In de cursus Integralen kon je de lengte van de parabool $y = x^2$ op $[0, 1]$ alleen numeriek berekenen: $\int_0^1 \sqrt{1 + 4x^2}\dx \approx 1{,}479$. Lukt het nu exact? Met $x = \tfrac12\tan u$ krijg je $\int \tfrac{1}{\cos^3 u}\du$, en dat is veel werk. Korter is een truc uit Integralen: partiële integratie, tot de integraal terugkeert.

\begin{theorie}
\begin{stelling}{Primitieve van $\sqrt{x^2 + a^2}$}
Voor $a > 0$ is \qquad $\formule{\int \sqrt{x^2 + a^2}\dx = \dfrac12\Bigl(x\sqrt{x^2 + a^2} + a^2\ln\bigl(x + \sqrt{x^2 + a^2}\,\bigr)\Bigr) + C}$
\end{stelling}

\begin{bewijs}
Noem de integraal $I$, en neem $f = \sqrt{x^2 + a^2}$ en $g' = 1$, dus $f' = \dfrac{x}{\sqrt{x^2 + a^2}}$ en $g = x$. Partiële integratie geeft
\bewijsregel{I = x\sqrt{x^2 + a^2} - \ig \dfrac{x^2}{\sqrt{x^2 + a^2}}\dx = x\sqrt{x^2 + a^2} - \ig \dfrac{(x^2 + a^2) - a^2}{\sqrt{x^2 + a^2}}\dx .}
De laatste integraal is $I - a^2\ln\bigl(x + \sqrt{x^2 + a^2}\bigr)$, volgens de vorige stelling. Dus
\bewijsregel{2I = x\sqrt{x^2 + a^2} + a^2\ln\bigl(x + \sqrt{x^2 + a^2}\,\bigr) + C .}
\end{bewijs}
\end{theorie}

\voorbeeld{}{Bereken de lengte van de parabool $y = x^2$ op $[0, 1]$ exact.\label{vb:h2-parabool}}{%
\uitleg{De formule voor de booglengte geeft $\int_0^1 \sqrt{1 + 4x^2}\dx$. Haal de factor 4 uit de wortel, dan staat er de vorm van de stelling, met $a = \tfrac12$.}
\naast{\begin{stappenvrij}{7.6cm}{3.3cm}
\sqrt{1 + 4x^2} = 2\sqrt{x^2 + \tfrac14} & $a^2 = \tfrac14$. \\[1ex]
L = \grens{x\sqrt{x^2 + \tfrac14} + \tfrac14\ln\bigl(x + \sqrt{x^2 + \tfrac14}\bigr)}{0}{1} & De factor 2 valt weg tegen de $\tfrac12$ uit de formule. \\[3ex]
= \dfrac{\sqrt5}{2} + \dfrac14\ln\Bigl(1 + \dfrac{\sqrt5}{2}\Bigr) - \dfrac14\ln\dfrac12 & $\sqrt{1 + \tfrac14} = \tfrac{\sqrt5}{2}$. \\[3ex]
= \dfrac{\sqrt5}{2} + \dfrac14\ln\bigl(2 + \sqrt5\bigr) \approx 1{,}4789 & Twee logaritmen samen.
\end{stappenvrij}}{\begin{tikzpicture}[x=1.8cm,y=1.8cm,font=\scriptsize]
  \rooster{0}{1}{1}{0}{1}{1}
  \kromme[very thick]{0}{1}{(\x)^2}
  \draw[accent,dashed] (0,0) -- (1,1);
\end{tikzpicture}}
Dat is de numerieke waarde uit Integralen, nu exact.}

\begin{werkwijze}
Bij een wortel van $a^2 + x^2$: stel $x = a\tan u$ met $u \in \bigl]-\tfrac{\pi}{2}, \tfrac{\pi}{2}\bigr[$, dan is $\sqrt{a^2 + x^2} = \tfrac{a}{\cos u}$ en $\mathrm{d}x = \tfrac{a}{\cos^2 u}\du$. Om terug te keren naar $x$ teken je een driehoek met overstaande zijde $x$, aanliggende zijde $a$ en schuine zijde $\sqrt{a^2 + x^2}$. Voor $\tfrac{1}{\sqrt{x^2 + a^2}}$ en $\sqrt{x^2 + a^2}$ gebruik je meteen de formules.
\end{werkwijze}

\nujij{Bereken $\int_0^{4/3} \dfrac{1}{\sqrt{x^2 + 1}}\dx$.}{$\ln\bigl(\tfrac43 + \tfrac53\bigr) - \ln 1 = \ln 3$.}

\begin{oefeningen}
\oef{h2-tan}
\end{oefeningen}

%========================================================================================
\section{Een wortel van \texorpdfstring{$x^2 - a^2$}{x²-a²}}

Blijft $x^2 - a^2$. Nu moet een kwadraat min 1 een kwadraat geven. Lees dezelfde formule anders: $\dfrac{1}{\cos^2 u} - 1 = \tan^2 u$. Stel dus $x = \dfrac{a}{\cos u}$. Voor $x \geqslant a$ neem je $u \in \bigl[0, \tfrac{\pi}{2}\bigr[$, en dan is $\sqrt{x^2 - a^2} = a\tan u$ en $\mathrm{d}x = \dfrac{a\sin u}{\cos^2 u}\du$.

\voorbeeld{}{Bereken $\ig_2^4 \dfrac{\sqrt{x^2 - 4}}{x}\dx$.\label{vb:h2-sec}}{%
\uitleg{Stel $x = \tfrac{2}{\cos u}$. De grenzen: $x = 2$ geeft $\cos u = 1$, dus $u = 0$, en $x = 4$ geeft $\cos u = \tfrac12$, dus $u = \tfrac{\pi}{3}$.}
\begin{stappenbreed}
\text{Stel } x = \dfrac{2}{\cos u}, \quad \mathrm{d}x = \dfrac{2\sin u}{\cos^2 u}\du & De afgeleide van $(\cos u)^{-1}$ is $\tfrac{\sin u}{\cos^2 u}$. \\[2.5ex]
\ig_0^{\pi/3} \dfrac{2\tan u}{2/\cos u} \cdot \dfrac{2\sin u}{\cos^2 u}\du = \ig_0^{\pi/3} 2\tan^2 u\du & $\tan u \cdot \cos u = \sin u$, en $\tfrac{\sin^2 u}{\cos^2 u} = \tan^2 u$. \\[3ex]
= 2\grens{\tan u - u}{0}{\pi/3} = 2\sqrt3 - \dfrac{2\pi}{3} \approx 1{,}37 & Zoals in Integralen: $\tan^2 u = \tfrac{1}{\cos^2 u} - 1$.
\end{stappenbreed}}

\begin{werkwijze}
Bij een wortel van $x^2 - a^2$ met $x \geqslant a$: stel $x = \tfrac{a}{\cos u}$ met $u \in \bigl[0, \tfrac{\pi}{2}\bigr[$, dan is $\sqrt{x^2 - a^2} = a\tan u$ en $\mathrm{d}x = \tfrac{a\sin u}{\cos^2 u}\du$. Terug naar $x$: $\cos u = \tfrac ax$.
\end{werkwijze}

\nujij{Bereken $\int_{\sqrt2}^2 \dfrac{1}{x\sqrt{x^2 - 1}}\dx$.}{$x = \tfrac{1}{\cos u}$: de integrand wordt 1, met grenzen $\tfrac{\pi}{4}$ en $\tfrac{\pi}{3}$. Uitkomst $\tfrac{\pi}{12}$.}

\begin{oefeningen}
\oef{h2-sec}
\end{oefeningen}

%========================================================================================
\section{De t-formules}

\subsection{Een substitutie die de goniometrie wegwerkt}

Bij $\int \dfrac{1}{3 + 5\cos x}\dx$ helpt geen van de vorige substituties: er staat geen wortel, en ook geen factor $\sin x$ die de afgeleide van $\cos x$ zou zijn. Kunnen we de sinus en de cosinus helemaal kwijtraken? Ja: allebei zijn ze een rationale functie van één getal, de tangens van de halve hoek. De figuur toont waarom. Trek vanuit $A(-1, 0)$ een rechte naar het punt $P(\cos x, \sin x)$ van de eenheidscirkel. Volgens de stelling van de omtrekshoek maakt die rechte een hoek $\tfrac x2$ met de $x$-as, dus haar richtingscoëfficiënt is $t = \tan\tfrac x2$. Omgekeerd ligt $P$ vast door die rechte: het tweede snijpunt met de cirkel.

\begin{theorie}
\begin{minipage}[t]{.64\linewidth}
\vspace{0pt}
\begin{stelling}{t-formules}
Is $x \in \,]-\pi, \pi[$ en $t = \tan\dfrac x2$, dan is

\vspace{1mm}
\qquad $\formule{\sin x = \dfrac{2t}{1 + t^2}}$ \qquad en \qquad $\formule{\cos x = \dfrac{1 - t^2}{1 + t^2}}$
\end{stelling}
\end{minipage}\hfill
\begin{minipage}[t]{.32\linewidth}
\vspace{0pt}
\centering
\begin{tikzpicture}[scale=1.3,font=\scriptsize]
  \draw[->,grijs] (-1.3,0) -- (1.35,0);
  \draw[->,grijs] (0,-.4) -- (0,1.3);
  \draw[accent,thick] (0,0) circle (1);
  \draw[fel!80!black,thick] (-1,0) -- (.5,.866);
  \draw[lijngrijs] (0,0) -- (.5,.866);
  \draw (-.65,0) arc[start angle=0,end angle=30,radius=.35];
  \node at (-.48,.1) {$\frac x2$};
  \draw (.25,0) arc[start angle=0,end angle=60,radius=.25];
  \node at (.34,.2) {$x$};
  \fill (-1,0) circle (1.2pt) node[below left,inner sep=1pt] {$A$};
  \fill[accent] (.5,.866) circle (1.4pt) node[above right,inner sep=1pt,text=black] {$P$};
  \fill (0,.577) circle (1.2pt) node[left,inner sep=2pt] {$t$};
\end{tikzpicture}
\end{minipage}

\begin{bewijs}
Met de formules voor de dubbele hoek, en $\cos^2\tfrac x2 + \sin^2\tfrac x2 = 1$:
\bewijsregel{\sin x = \dfrac{2\sin\frac x2\cos\frac x2}{\cos^2\frac x2 + \sin^2\frac x2}, \qquad \cos x = \dfrac{\cos^2\frac x2 - \sin^2\frac x2}{\cos^2\frac x2 + \sin^2\frac x2} .}
Deel telkens teller en noemer door $\cos^2\tfrac x2 \neq 0$, want $\tfrac x2 \in \bigl]-\tfrac{\pi}{2}, \tfrac{\pi}{2}\bigr[$. Dan staat er $\dfrac{2t}{1 + t^2}$ en $\dfrac{1 - t^2}{1 + t^2}$.
\end{bewijs}
\end{theorie}

Om te substitueren heb je ook $\mathrm{d}x$ nodig. Uit $t = \tan\tfrac x2$ volgt $x = 2\bgtan t$, en dus $\mathrm{d}x = \dfrac{2}{1 + t^2}\dt$. Een rationale functie van $\sin x$ en $\cos x$ wordt zo een rationale functie van $t$, en die kun je met hoofdstuk 1 altijd integreren.

\voorbeeld{}{Bereken $\ig \dfrac{1}{3 + 5\cos x}\dx$.\label{vb:h2-t}}{%
\uitleg{Een breuk met $\cos x$ zonder bruikbare afgeleide: de t-formules. Daarna staat er een rationale functie, met partieelbreuken.}
\begin{stappenbreed}
3 + 5\cos x = \dfrac{3(1 + t^2) + 5(1 - t^2)}{1 + t^2} = \dfrac{8 - 2t^2}{1 + t^2} & Stel $t = \tan\tfrac x2$, op gelijke noemer. \\[3ex]
\ig \dfrac{1 + t^2}{8 - 2t^2} \cdot \dfrac{2}{1 + t^2}\dt = \ig \dfrac{1}{4 - t^2}\dt & De factor $1 + t^2$ valt weg, zoals bijna altijd. \\[3ex]
\dfrac{1}{(2 - t)(2 + t)} = \dfrac14\Bigl(\dfrac{1}{2 - t} + \dfrac{1}{2 + t}\Bigr) & Partieelbreuken: bedek en vul $t = \pm 2$ in. \\[3ex]
= \dfrac14\ln\Bigl|\dfrac{2 + t}{2 - t}\Bigr| + C = \dfrac14\ln\Bigl|\dfrac{2 + \tan\frac x2}{2 - \tan\frac x2}\Bigr| + C & Terug naar $x$.
\end{stappenbreed}}

\voorbeeld{}{Bereken $\ig_0^{\pi/2} \dfrac{1}{5 + 4\cos x}\dx$.\label{vb:h2-t2}}{%
\uitleg{Met grenzen gaan die meteen mee: $x = 0$ geeft $t = 0$, en $x = \tfrac{\pi}{2}$ geeft $t = \tan\tfrac{\pi}{4} = 1$.}
\naast{\begin{stappenvrij}{6.6cm}{3.4cm}
5 + 4\cos x = \dfrac{9 + t^2}{1 + t^2} & $5(1 + t^2) + 4(1 - t^2)$. \\[2.5ex]
\ig_0^1 \dfrac{2}{9 + t^2}\dt = \dfrac23\grens{\bgtan\dfrac t3}{0}{1} & Nu een noemer zonder nulwaarden. \\[3ex]
= \dfrac23\bgtan\dfrac13 \approx 0{,}215 & Controle: de integrand ligt tussen $\tfrac19$ en $\tfrac15$, op een interval met breedte $1{,}57$.
\end{stappenvrij}}{\begin{tikzpicture}[x=1.4cm,y=7cm,font=\scriptsize]
  \gebied{0}{1.5708}{1/(5+4*cos(\x r))}
  \draw[black!14,very thin,xstep=.5,ystep=.1] (0,0) grid (2,.3);
  \draw[->,grijs] (0,0) -- (2.2,0) node[right,text=black] {$x$};
  \draw[->,grijs] (0,0) -- (0,.35) node[above,text=black] {$y$};
  \kromme{0}{2}{1/(5+4*cos(\x r))}
  \node[below] at (1.5708,0) {$\frac{\pi}{2}$};
  \node[left] at (0,.1) {0,1};
  \node[left] at (0,.2) {0,2};
\end{tikzpicture}}}

\begin{werkwijze}
Is de integrand een breuk met $\sin x$ en $\cos x$, en helpt geen eenvoudigere substitutie, stel dan $t = \tan\tfrac x2$. Vervang $\sin x$, $\cos x$ en $\mathrm{d}x$ door $\tfrac{2t}{1 + t^2}$, $\tfrac{1 - t^2}{1 + t^2}$ en $\tfrac{2}{1 + t^2}\dt$. Integreer de rationale functie van $t$, en reken de grenzen om met $t = \tan\tfrac x2$. Pas op: het interval moet binnen $]-\pi, \pi[$ liggen, want in $x = \pm\pi$ bestaat $t$ niet.
\end{werkwijze}

\nujij{Bereken $\int \dfrac{1}{\sin x}\dx$ op $]0, \pi[$.}{$\int \dfrac{1 + t^2}{2t} \cdot \dfrac{2}{1 + t^2}\dt = \ln|t| + C = \ln\tan\dfrac x2 + C$.}

\begin{oefeningen}
\oef{h2-t}
\end{oefeningen}

\subsection{Niet altijd de kortste weg}

De t-formules werken altijd, maar het rekenwerk loopt snel op. Kijk dus eerst of het eenvoudiger kan. Bij $\int \dfrac{\sin x}{2 + \cos x}\dx$ staat de afgeleide van de noemer op het teken na in de teller: dat geeft $-\ln(2 + \cos x) + C$, zonder t-formules. En bij een oneven macht van $\sin x$ of $\cos x$ hou je één factor apart en substitueer je de andere functie, zoals in Integralen.

\begin{oefeningen}
\oef{h2-kort}
\end{oefeningen}

%========================================================================================
\section{Welke substitutie?}

De substitutie volgt uit de vorm van de integrand. Toch lijken integralen die heel anders aangepakt worden vaak op elkaar, en soms is er helemaal geen goniometrie nodig.

\voorbeeld{}{Vijf integralen met $1 - x^2$. (Niet uitgewerkt, controleer met de afgeleide.)\label{vb:h2-keuze}}{%
\uitleg{Kijk naar de bouw: een wortel of niet, en een factor $x$ in de teller of niet.}
\begin{stappenbreed}
\ig \dfrac{x}{\sqrt{1 - x^2}}\dx = -\sqrt{1 - x^2} + C & Substitutie $u = 1 - x^2$: de $x$ is op een factor na de afgeleide. \\[3ex]
\ig \dfrac{1}{\sqrt{1 - x^2}}\dx = \bgsin x + C & Een standaardvorm. \\[3ex]
\ig \sqrt{1 - x^2}\dx = \dfrac12\bgsin x + \dfrac x2\sqrt{1 - x^2} + C & Stel $x = \sin u$, of de formule. \\[3ex]
\ig \dfrac{1}{1 - x^2}\dx = \dfrac12\ln\Bigl|\dfrac{1 + x}{1 - x}\Bigr| + C & Geen wortel: partieelbreuken. \\[3ex]
\ig \dfrac{1}{(1 - x^2)^{3/2}}\dx = \dfrac{x}{\sqrt{1 - x^2}} + C & Stel $x = \sin u$.
\end{stappenbreed}}

Staat er een wortel van een veelterm van de tweede graad met een term in $x$, dan splits je eerst het kwadraat af. Dan zie je welke van de drie vormen het is.

\voorbeeld{}{Bereken $\ig \dfrac{1}{\sqrt{3 + 2x - x^2}}\dx$.\label{vb:h2-afsplits}}{%
\uitleg{Splits het kwadraat af: er staat een $-x^2$, dus het wordt een vorm $a^2 - (x - m)^2$.}
\begin{stappen}
3 + 2x - x^2 = 4 - (x - 1)^2 & $-(x^2 - 2x + 1) + 1 + 3$. \\[1ex]
\ig \dfrac{1}{\sqrt{2^2 - (x - 1)^2}}\dx = \bgsin\dfrac{x - 1}{2} + C & De standaardvorm met $u = x - 1$ en $a = 2$.
\end{stappen}}

\begin{werkwijze}
\begin{tabular}{@{}>{$}l<{$}@{\hspace{5mm}}>{$}l<{$}@{\hspace{5mm}}>{$}l<{$}@{}}
\text{wortel van } a^2 - x^2 & \text{stel } x = a\sin u & \text{dan } \sqrt{a^2 - x^2} = a\cos u \\[.5ex]
\text{wortel van } a^2 + x^2 & \text{stel } x = a\tan u & \text{dan } \sqrt{a^2 + x^2} = \tfrac{a}{\cos u} \\[.5ex]
\text{wortel van } x^2 - a^2 & \text{stel } x = \tfrac{a}{\cos u} & \text{dan } \sqrt{x^2 - a^2} = a\tan u \\[.5ex]
\text{breuk met } \sin x \text{ en } \cos x & \text{stel } t = \tan\tfrac x2 & \text{dan een rationale functie van } t
\end{tabular}

\vspace{1mm}
Kijk eerst of een gewone substitutie of een standaardvorm volstaat. Splits het kwadraat af als er een term in $x$ staat.
\end{werkwijze}

\begin{oefeningen}
\oef{h2-keuze}
\end{oefeningen}

%========================================================================================
\opdrachtenkop

\begin{oefeningen}
\needspace{18\baselineskip}
\oef{h2-o1}

\oef{h2-o2}

\oef{h2-o3}

\oef{h2-o4}

\oef{h2-o5}

\oef{h2-o6}

\oef{h2-o7}

\oef{h2-o8}

\oef{h2-o9}

\oef{h2-o10}

\oef{h2-o11}

\oef{h2-o12}

\oef{h2-o13}

\oef{h2-o14}

\oef{h2-o15}

\oef{h2-o16}

\oef{h2-o17}

\oef{h2-o18}

\oef{h2-o19}

\oef{h2-o20}
\end{oefeningen}

%----------------------------------------------------------------------------------------
% KOPPEL: verbanden met andere leerstof
%----------------------------------------------------------------------------------------
\needspace{20\baselineskip}
\opdrachtsectie{Koppel}

\lkr{Goniometrische functies en Complexe getallen zijn gezien in het 5de jaar. Oppervlakken is een module van het leerplan Complementair: de oefening geeft de vergelijking van de torus, en kan ook zonder die module.}

\begin{oefeningen}
\oef{h2-koppel-1}

\oef{h2-koppel-2}

\oef{h2-koppel-3}
\end{oefeningen}

%----------------------------------------------------------------------------------------
\needspace{30\baselineskip} % de test en het schema Fout? Herhaal samen op een pagina
\opdrachtsectie{Test jezelf}

\begin{oefeningen}
\oef{h2-t1}

\oef{h2-t2}

\oef{h2-t3}

\oef{h2-t4}

\oef{h2-t5}

\oef{h2-t6}

\oef{h2-t7}
\end{oefeningen}

\herhaalschema{%
  \fv{oef:h2-t1}{vb:h2-sin}{oef:h2-sin}%
  \fv{oef:h2-t2}{vb:h2-tan}{oef:h2-tan}%
  \fv{oef:h2-t3}{vb:h2-parabool}{oef:h2-tan}%
  \fv{oef:h2-t4}{vb:h2-sec}{oef:h2-sec}%
  \fv{oef:h2-t5}{vb:h2-t2}{oef:h2-t}%
  \fv{oef:h2-t6}{vb:h2-terug}{oef:h2-sin}%
  \fv{oef:h2-t7}{vb:h2-keuze}{oef:h2-keuze}}

\ifzelfstudie
\section*{Maak je eigen samenvatting}
Sluit dit boek en schrijf in je schrift: de substitutie $x = g(u)$ met het bewijs en wat er met de grenzen gebeurt, de drie goniometrische substituties met de driehoek voor de terugweg, de primitieven van $\sqrt{a^2 - x^2}$, $\tfrac{1}{\cos x}$, $\tfrac{1}{\sqrt{x^2 + a^2}}$ en $\sqrt{x^2 + a^2}$ met het idee van hun bewijs, de t-formules met de figuur en $\mathrm{d}x$, en hoe je een substitutie kiest. Vergelijk daarna met de kaders.
\fi
